\documentclass[12pt,a4paper,oneside]{article}
\usepackage[brazil]{babel}
\usepackage[utf8]{inputenc}
\usepackage{olymp}

\setcounter{problem}{2}

\begin{document}

\begin{problem}{Número perfeito}{perfeito}{2}
 
Um número perfeito é um número natural cuja soma de seus divisores, exceto ele mesmo, é igual ao próprio número. Por exemplo, 6 é um número perfeito, pois $1+2+3=6$.

Faça um programa que, dado um número natural $N$, determine se ele é ou não um número perfeito.

%\Notes
%
%Observações, se tiver.

\InputFile

A entrada contém apenas um número natural $N$. Restrição: $1 \leq N \leq 10000$.


\OutputFile

Seu programa deve gerar apenas uma linha na saída, contendo a palavra ``\texttt{SIM}'' caso $N$ seja um número perfeito ou a palavra ``\texttt{NAO}'' caso contrário.

% \Notes
% Note que a mensagem escrita não tem acento e tem um ponto final.

\Example

\begin{example}
\exmp{
6
}{
SIM
}%
\end{example}

\begin{example}
\exmp{
28
}{
SIM
}%
\end{example}

\begin{example}
\exmp{
360
}{
NAO
}%
\end{example}



%Comentários sobre o exemplo, se necessário.

\end{problem}

\end{document}
